\chapter{PHÂN TÍCH VÀ THẢO LUẬN}
\label{chap:phantich}

\section{VAI TRÒ CỦA TỪNG THÀNH PHẦN TRONG MCTS}
\begin{figure}[htbp]
  \includegraphics[width=\linewidth]{H5_mcts_stages.pdf}
  \caption{Pass@1 của RethinkMCTS qua từng bước: chương trình đầu tiên, sau search, sau critic}
  \label{fig:buoc}
\end{figure}
\lipsumtodo{Search so với chương trình đầu tiên; critic chọn lại offline (B5.2).}

\section{HÀNH VI TÌM KIẾM}
\lipsumtodo{Độ sâu cây, số ứng viên, nguồn gốc lời giải đúng, dừng sớm tính lại offline (B5.4).}

\section{BÀI HỌC TỪ QUÁ TRÌNH TÁI HIỆN}
\lipsumtodo{Các lệch của cài đặt cũ và cách v3 sửa (B5.3, D-006); mô tả, không dùng số liệu v1.}

\section{PHÂN LOẠI LỖI VÀ NGHIÊN CỨU TÌNH HUỐNG}
\lipsumtodo{Chỉ chọn case từ run v3.}

\section{CÁC YẾU TỐ ẢNH HƯỞNG ĐẾN TÍNH HỢP LỆ}
\label{sec:hople}
\lipsumtodo{Một seed; nhiễm dữ liệu; temperature 0 trên vLLM không tất định tuyệt đối; Qwen3.5-9B bị cắt output trên APPS và chạy trên hai nền tảng (D-010); Competition có ít hidden test.}
